\chapter{绪论}

\section{课题背景及意义}

本节填写课程设计的背景与意义。可从相关技术的发展现状、工程应用需求、课程设计训练目标等方面展开说明，说明为什么选择该课题、该课题解决什么问题，以及完成该课题对专业学习和工程实践能力培养的意义。

撰写时可先介绍课题所属领域的基本背景，再说明现有方案或传统方法中存在的问题，最后引出本课程设计的研究对象和设计价值。文字应避免直接堆砌概念，重点说明本课题与课程内容、硬件平台、软件设计或工程实现之间的关系。

本节也可以结合具体课程要求，说明该课题涉及的主要知识点，例如系统总体设计、硬件资源分析、软件模块划分、程序调试、功能测试和结果分析等。通过本课题的设计与实现，可以加深对相关理论知识和工程开发流程的理解。

\section{课题任务及目标}

本节填写课程设计的主要任务和设计目标。可根据课程设计任务书，说明本课题需要完成的基本功能、扩展功能和验证要求。对于未完成或不属于本设计范围的内容，应在此处进行说明，避免正文中出现功能范围不清的问题。

本课程设计的主要任务可以按照以下逻辑描述：首先明确系统使用的硬件平台或软件环境；其次说明系统需要实现的核心功能；然后说明各功能模块之间的配合关系；最后说明系统需要达到的测试效果和验收目标。

本课程设计的目标包括：完成系统总体方案设计，明确主要硬件资源和软件模块；完成关键功能模块的设计与实现；完成系统调试与测试，并通过实验现象或测试结果验证系统功能的正确性和稳定性。具体目标可根据实际课题内容进一步补充和修改。

\section{论文结构安排}

本文围绕课程设计任务展开，全文结构安排如下。

第一章为绪论，主要介绍课题背景、研究意义、设计任务和设计目标。

第二章为系统总体方案设计，主要说明系统需求分析、总体结构、功能模块划分和整体工作流程。

第三章为系统硬件设计，主要介绍系统所用硬件平台、主要器件、接口连接关系和硬件资源分配。

第四章为系统软件设计，主要说明软件总体结构、关键模块设计、程序流程和核心实现方法。

第五章为系统测试与结果分析，主要介绍测试环境、测试过程、实验结果和存在的问题，并对课程设计完成情况进行总结。

\section{本章小结}

本章主要介绍了课程设计的课题背景、设计意义、任务范围和总体目标，并对全文结构进行了简要说明。通过本章内容，可以明确本课程设计需要解决的主要问题和后续章节的展开思路。